%% \cref on an example.  Without the declarations linguexx makes, cleveref
%% has no name for the example counters and prints its unknown-type marker:
%% "?? (1)".  The names are declared EMPTY on purpose -- linguistics prose
%% refers to an example by its number alone, and \theExNo already brings the
%% parentheses -- so \cref prints what \ref prints, and what it contributes
%% is its list and range handling.
\input{_preamble}
\usepackage{cleveref}
\begin{document}
\ex.\label{cv:one} CVONE first.
\a.\label{cv:a}\sublabel{cv:sa} CVA letter a.
\b.\label{cv:b}\sublabel{cv:sb} CVB letter b.
\a.\label{cv:r} CVR roman i.
\z.\z.

\ex.\label{cv:two} CVTWO second.

CVFN host.\footnote{\ex.\label{cv:fn} CVFNEX in a footnote.

}

CVSINGLE \cref{cv:one} CVSUB \cref{cv:a} CVROMAN \cref{cv:r}
CVMULTI \cref{cv:one,cv:two} CVRANGE \cref{cv:a,cv:b}
CVCAP \Cref{cv:one} CVFNREF \cref{cv:fn} CVPLAIN \ref{cv:a}

%% \crefrange is cleveref's OTHER combining form, and it goes through the
%% empty \crefname declarations by a different route than \cref does (the
%% range conjunction, not the list one), so \cref lists alone did not cover
%% it: a \crefname guard that let a name through would show up here as
%% "example (1) to example (2)".
%% Each of the two lines below is its own paragraph, short enough not to
%% reach the right margin: a sentinel hyphenated across a line break would
%% make these assertions fail for a reason that has nothing to do with
%% cleveref.
CVCREFRANGE \crefrange{cv:a}{cv:b} CVCREFRANGEMAIN \crefrange{cv:one}{cv:two}

%% ... and the reason linguexx's own \refrange is NOT redundant with it:
%% cleveref spells both endpoints out in full, while \refrange compresses a
%% shared prefix to "(1a--b)".  The two must therefore print DIFFERENTLY
%% over the same pair; nothing asserted that before, so "cleveref already
%% does ranges" looked like grounds for deleting \refrange.
CVREFRANGE \refrange{cv:sa}{cv:sb}
\end{document}
